% A chapter fragment. The main document supplies theorem environments and
% the leanblueprint macros \lean, \leanok, and \uses.
% Only the two HeaderKernel nodes below are marked leanok.

\section{A boundary-header repair of coloured composition}

This chapter proves the hereditary-side core composition statement. The
flexible zero/rank definition requires hereditary listed sides, as defined
in Definition~\ref{def:hereditary}: (G3) allows a fine outer side to cross
a dummy parent endpoint, and for such a side the clause ``$b\subset\sigma'd$''
of Definition~\ref{def:control/substitution} is undefined. Accordingly
\emph{every theorem and lemma of this chapter carries the standing
hypothesis that the macrotemplate system has hereditary listed sides}, so
that each fine outer side lies in a unique scaled parent listed side; the
hypothesis is repeated in the statements of Theorems~\ref{thm:compose},
\ref{lem:lift-hierarchy} and~\ref{thm:comp/soundness} and in their Lean
targets. Original nonhereditary dummy cuts are
handled separately by the strict schedule proof in the final chapter.
This chapter proves a corrected composition statement for an already
constructed macrotemplate system. It is conditional on the finite geometric,
routing, and separation certificate of Definition~\ref{def:certificate}.
It does not establish that certificate for every STS substitution. In
particular, the arbitrary assignment of tiles to grid squares in the source
article does not establish connectedness or the required routing data.
Every internal interface has a distinct finite ID naming its exact two
owner-side slots and relative displacement. Identity data are triples
$(s,r,e)$; fingerprints project them to the ordered type pair $(s,r)$,
so all occurrence counts are unchanged.
All declarations in the namespace \texttt{Tilings.Planned} below are proposed
formalization targets, with mathematical proofs given here; they have not
been compiled. Only the two explicitly marked nodes in
\texttt{Tilings.HeaderKernel} refer to the compiled Lean kernel.

The certificate interfaces used here must include the following conventions.
Outer sides are marked. A child's type determines its owner template and
therefore the form of its parent. Print maps depend only on that form, and
copy-source maps depend only on the child position and the parent form.
The boundary belt is connected, the central tile has no outer or boundary
side, and identity data determine rigid relative placements. All polygon
patches have disjoint interiors and cover their parent polygon. These
conventions are essential uses of the certificate, rather than conclusions
derived from normality alone.

\begin{definition}[Typed decorations and normality]
\label{def:control/typed-decoration}
\lean{Tilings.Planned.typedDecoration}
\uses{def:certificate,def:substitution}
Let $F$ be the finite set of forms and $T$ the finite set of positioned child
types. Write $\operatorname{form}:T\to F$, and let $S(t)$ be the finite
indexed side set of $\operatorname{form}(t)$. The owner of a child type is
its parent form, not necessarily its own form. Side identifiers are retained
even when consecutive polygonal sides are collinear.

An index belongs to a finite tagged alphabet $U$. Its tags distinguish zero,
positive alignment ranks, type labels, exact interface identities, the symbol
$\bot$, and the headers defined below. A decoration of type $t$ consists
of a five-tuple $c_A(d)\in U^5$ on each $d\in S(t)$; its centre stores $t$.
In the hierarchical variant considered here, position $2$ is zero. The
geometric polygon has form $\operatorname{form}(t)$.

The local cells are all cells of an unmarked side, the fifth cell of a special
side, and the first and fifth cells of a marked ordinary side. A decoration
is \emph{normal} when these cells have the values prescribed by its type.
Thus ordinary inner first cells contain the relevant ordered interface identity.
On a special side, the local fifth value is the interface identity if that side
is a print in its owner template, and zero if it is not a print. In particular,
special does not mean print. Call a decoration \emph{identity-correct} when
every inner nonprint side carries its prescribed identity in position $1$
and every print carries it in position $5$. Normality alone does not impose
the first value on a special nonprint side. In the flexible variant of
this chapter an outer fourth cell is not local: an outer side is marked,
and its zero-versus-rank value depends on the parent's mode. In the strict
variant of the final chapter the outer fourth cell is the generated
schedule $W$, which is determined by the type alone and is therefore local
there. The list of local cells relies on the marking discipline (G7): an
unmarked side is inner, nonprint, nonboundary and on no path, so all its
cells are fixed by the type.
\end{definition}

\begin{lemma}[A tile with an outer side has a boundary side]
\label{lem:comp/outer-boundary}
\lean{Tilings.Planned.boundary_side_of_outer_side,Tilings.Planned.Strict.exists_boundary_side_of_outer}
\uses{def:certificate}
In a macrotemplate every child type with an outer side also has a boundary
side, that is, an inner side with an endpoint on the macroboundary.
\end{lemma}
\begin{proof}
By (G1) the type has an inner side, so not all of its sides are outer.
Walk around its boundary polygon from an outer side; the first side that
is not outer starts at the endpoint of the last outer side, which lies on
the macroboundary. That side is inner (the template covers the parent
polygon, so a side not on the macroboundary is shared) and has an endpoint
on the macroboundary. Collinear listed vertices do not affect the walk.
\end{proof}

\begin{definition}[The finite parent header]
\label{def:control/header}
\lean{Tilings.HeaderKernel.makeHeader}
\leanok
\uses{def:control/typed-decoration}
For a type $t$, a Boolean outer-side predicate $o_t:S(t)\to\{0,1\}$,
and a fourth-index vector $a:S(t)\to U$, set
\[
 z_{t,a}(d)=o_t(d)\mathbin{\wedge}[a(d)=0],\qquad
 H(t,a)=\bigl\langle t,z_{t,a}\bigr\rangle
 \in\sum_{t\in T}\bigl(S(t)\to\{0,1\}\bigr).
\]
The Lean definitions are \texttt{Header}, \texttt{zeroMode}, and
\texttt{makeHeader}, with precisely these dependent types and formulas.
The sum is finite when $T$ and each $S(t)$ are finite; the compiled instance
\texttt{headerFintype} uses classical decidable equality and is
\texttt{noncomputable}. For a decorated tile $A$ write $H(A)$ for the header
computed from its type and its fourth vector. This is computed data, not a
new local field fixed by the type.
\end{definition}

\begin{lemma}[Recovery of a header from fourth-port data]
\label{lem:comp/port-recovery}
\lean{Tilings.HeaderKernel.makeHeader_eq_of_port_recovery}
\leanok
\uses{def:control/header}
Fix a type $t$ and fourth vectors $a,b:S(t)\to U$. If a vector
$p:S(t)\to U$ satisfies $p(d)=a(d)$ and $p(d)=b(d)$ for every $d$, then
$H(t,a)=H(t,b)$.
\end{lemma}
\begin{proof}
 \leanok
For each $d$, transitivity gives $a(d)=p(d)=b(d)$. Consequently
$o_t(d)\wedge[a(d)=0]=o_t(d)\wedge[b(d)=0]$. Function extensionality
identifies the two Boolean vectors, and their type components are already
the same. This is the proof implemented in the kernel. If types are initially
identified by an equality $t=s$, their dependent side sets must first be
transported along that equality; the compiled auxiliary lemma
\texttt{makeHeader\_eq\_of\_aligned\_fourth} performs this transport.
\end{proof}

\begin{definition}[The repaired decorated substitution]
\label{def:control/substitution}
\lean{Tilings.Planned.headerSubstitution}
\uses{def:certificate,def:substitution,def:control/header}
The uncoloured substitution $\sigma'$ and its finite templates are unchanged.
For a decorated parent $A$, its repaired decomposition $\widehat\sigma A$
has the prescribed child types. Prints carry the first four cells of the
corresponding parent side, and their fifth cell is the interface identity of the
shared interior side. This prescription applies to \emph{both} owners of a
print. A side which is not a print, including a special nonprint side,
has its prescribed first identity cell, zero second
cell, zero fifth cell, and third cell equal to $A_i(d)$ when it belongs to
the path $P_i(d)$, and zero otherwise. Its fourth cell is
\[
 B_4(b)=
 \begin{cases}
  0,&b\subset\sigma'd\text{ is outer and }z_A(d)=1,\\
  \operatorname{rank}(b),&b\text{ is outer and }z_A(d)=0,\\
  \operatorname{header}(H(A)),&b\text{ is boundary},\\
  \bot,&b\text{ is any remaining inner nonprint side}.
 \end{cases}
\]
The print prescription has precedence. Prints lie strictly inside the
macrotemplate and hence do not conflict with the boundary clause. Each
header tag is distinct from zero. The change from the article is the
boundary clause: it carries the entire parent header, rather than only the
parent type. This introduces finite control data and no new paths.
\end{definition}

\begin{definition}[Copy sources and origins]
\label{def:control/origin}
\lean{Tilings.Planned.copySource}
\uses{def:control/substitution,def:certificate}
Copied cells are positions $1$--$4$ on prints and position $3$ on ordinary
sides assigned to paths. Their source is the corresponding parent cell.
The source is determined by the child position and parent form. All other
cells, including the centre type and the synthesized boundary header, are
origins. Thus a boundary header is \emph{not} falsely treated as a literal
copy of the sixth cell.

Along an ancestor chain, repeatedly follow the source of a copied cell.
If a noncopied cell is eventually reached, it is the origin of the chain.
Otherwise the copied cell has no origin. A print or path may copy a
header-valued fourth cell; its source chain then stops at the synthesized
boundary cell. Values are unchanged along actual copy edges.
\end{definition}

\begin{definition}[Legality and properness]
\label{def:control/legal}
\lean{Tilings.Planned.headerLegal}
\uses{def:control/origin,def:control/substitution,def:tiling}
A decoration $A$ is legal when it has an infinite ancestor chain
$A=A_0,A_1,A_2,\ldots$, with $A_i$ a child of $A_{i+1}$, such that every
copied cell at every $A_i$ with no origin has value zero. All types and
colours lie in fixed finite sets; legality selects a subset of that finite
decorated tile set.

A side is nontrivial if its fourth and fifth cells are not both zero.
A plane tiling is proper if every nontrivial side has a complete, coincident
opposing side with the same full five-tuple. A trivial outer side may have
shifted contacts. Equality of colours is required only for contacts forced
by this rule.
\end{definition}

\begin{lemma}[Elementary legality properties]
\label{lem:comp/legal-properties}
\lean{Tilings.Planned.headerLegal_properties,Tilings.Planned.Strict.strictLegal_properties}
\uses{def:control/legal,def:control/typed-decoration,def:control/substitution}
A legal tile has a legal parent and is normal and identity-correct;
every child of a legal parent
is legal. All boundary fourth cells of a legal tile equal one parent header.
Third signals belonging to the same path are equal through a legal tile.
Every inner side of a legal tile is nontrivial. On an outer side, triviality
is equivalent to a zero fourth cell.
\end{lemma}
\begin{proof}
Delete the first term of a witnessing ancestor chain to obtain legality of
its parent. The child clauses fix every local cell, giving normality, and
supply first identities on all inner nonprint sides and fifth identities
on all prints, giving identity-correctness.
To prove legality of a child, prepend it to a chain for the parent. A copied
cell in the new first tile either reaches an origin there, or reaches the
corresponding parent cell. If the source chain is infinite, the parent's
legality forces its value, and hence the child's value, to be zero. All
requirements farther up the chain are unchanged.

The boundary and path equalities follow by choosing the first legal parent:
its single header supplies every boundary cell, and its single cell $A_i(d)$
supplies all appearances of a path signal in that child. At a print, the
fifth cell is a nonzero interface identity. At a boundary side, the fourth cell
is a nonzero header tag. At any other inner side, the fourth cell is $\bot$.
These cases exhaust inner sides. On outer sides the fifth cell is zero,
so the last assertion follows directly.
\end{proof}

\begin{lemma}[Assembly from active identity sides]
\label{lem:comp/assembly}
\lean{Tilings.Planned.assembleMacrotemplates,Tilings.Planned.Strict.strict_assembleMacrotemplates}
\uses{def:certificate,def:control/typed-decoration,def:control/legal,def:tiling}
A proper plane tiling of normal identity-correct typed tiles in which every
inner side is nontrivial admits a partition into positioned $\sigma'$ macrotemplates.
In particular, every proper legal tiling has such a partition.
\end{lemma}
\begin{proof}
On an inner nonprint side, identity-correctness gives its ordered identity
pair in the first cell. On a print, the fifth cell gives that pair.
Nontriviality forces a complete matching side. The pair identifies the
neighbouring type and the relative placement specified by the certificate.
Starting with one tile, traverse the connected finite inner-side graph of
its owner template. Each graph edge forces its required neighbour. Relative
placements agree because they are the placements of that fixed template;
an existing tile cannot occupy a different positive-area position while
matching the same rigid side placement.

Traversal therefore recovers the full template. Repeated traversal cannot
produce an additional occurrence at the same position, since tile interiors
are disjoint. Two recovered blocks sharing a tile have the same owner type
set and the same rigid placement, and hence are the same block. Every tile
belongs to a recovered block, so the blocks partition the plane. Trivial
outer contacts are not used in this argument. The final assertion follows
from Lemma~\ref{lem:comp/legal-properties}. Notice the hypothesis that inner
sides are active and identities are correct: normality alone would neither
force adjacency on an ordinary marked side whose fourth and fifth cells
were both zero, nor prescribe the first identity on a special nonprint side.
\end{proof}

\begin{lemma}[A common boundary header]
\label{lem:comp/boundary-header}
\lean{Tilings.Planned.boundary_header_constant,Tilings.Planned.Strict.strict_boundary_header_constant}
\uses{def:certificate,lem:comp/legal-properties,def:control/header}
Let $M$ be a proper finite coloured realization of the template $\sigma'\alpha$,
with the prescribed child types. Assume every tile except possibly the
centre is legal. Every boundary fourth cell then carries the same header
$h=(t,v)$, and $\operatorname{form}(t)=\alpha$.
\end{lemma}
\begin{proof}
Every boundary tile is noncentral and legal. Its boundary cells all equal
the header of any chosen legal parent. If two boundary tiles share a boundary
side, that side is nontrivial; properness identifies the two header tags.
Connectedness of the boundary belt propagates equality to every boundary
cell of $M$. Choose one boundary tile $B$ and a legal parent $C$ of it.
The type of $B$ determines the owner form $\alpha$, so $C$ has form $\alpha$.
Its header equals $h$; equality of type components gives $\operatorname{type}(C)=t$.
Thus $t$ has form $\alpha$, as required. Connectedness of this belt is an
explicit certificate obligation, not an inference from a sufficiently
large substitution scale.
\end{proof}

\begin{lemma}[Parent reconstruction and its exact header]
\label{lem:comp/reconstruction}
\lean{Tilings.Planned.reconstructed_header_eq,Tilings.Planned.Strict.strict_reconstructed_header_eq}
\uses{lem:comp/boundary-header,lem:comp/port-recovery,def:control/typed-decoration}
In the situation of Lemma~\ref{lem:comp/boundary-header}, define $D$ to have
type $t$ and, for every parent side $d$ and $1\le i\le5$, set
\[
 D_i(d)=\text{the third cell of port }i\text{ on }\sigma'd.
\]
Then $D$ is normal and identity-correct, $H(D)=h$, and every inner side of
$D$ is nontrivial.
\end{lemma}
\begin{proof}
Fix a local cell $(i,d)$ of $D$. Its port lies in a legal boundary tile $B$,
which has a boundary side by Lemma~\ref{lem:comp/outer-boundary} and hence
a legal parent $C$ with $H(C)=h$.
In particular $C$ has type $t$. The port clause says $D_i(d)=B_3(p_i)=C_i(d)$.
Since $C$ is normal of type $t$, this is the required local value of $D$.
Applying this separately to every local cell proves normality; no equality
of the whole parents $C$ chosen for different ports is needed.

For each $d$, choose its fourth port and a legal parent $C_d$ of that port
tile. We have $H(C_d)=h$ and $D_4(d)=(C_d)_4(d)$. Thus
\[
 v(d)=[d\text{ outer in }t]\wedge[(C_d)_4(d)=0]
     =[d\text{ outer in }t]\wedge[D_4(d)=0]=z_D(d).
\]
The dependent side sets have first been aligned using the established
equality of parent types. Extensionality gives $H(D)=h$. The kernel recovery
lemma is the algebraic component of this argument; obtaining the legal
parents and their common type remains part of this mathematical lemma.

For an inner nonprint side $d$, use its first port and a legal parent of
that port tile. Its type is $t$, so the child clauses force its first cell
to be the prescribed identity of $d$. This port recovers exactly $D_1(d)$.
For a print, the fifth cell of $D$ is its nonzero prescribed identity by
normality. Thus $D$ is identity-correct. For an inner nonprint side, the
legal parent $C_d$ chosen at its fourth port has a nonzero fourth cell: it
is a header or $\bot$. Hence $D_4(d)\ne0$. For a print, the normal fifth
identity cell is nonzero. Therefore every inner side of $D$ is nontrivial,
even before legality of $D$ has been established. This argument explicitly
covers special nonprints, whose fifth cell is zero.
\end{proof}

\begin{lemma}[Signals reach the prints]
\label{lem:comp/signals}
\lean{Tilings.Planned.path_and_print_recovery,Tilings.Planned.Strict.strict_path_and_print_recovery}
\uses{def:certificate,lem:comp/legal-properties,lem:comp/reconstruction}
Every nonprint third signal on $P_i(d)$ in $M$ equals $D_i(d)$.
The noncentral owner of the print of $d$ has first four cells
$(D_1(d),\ldots,D_4(d))$ on that print.
\end{lemma}
\begin{proof}
Begin at port $i$, whose third cell is $D_i(d)$ by construction.
Consecutive nonprint sides of the path belong to the same noncentral legal
tile, or are common sides of two such tiles. In the first case, their
signals agree because both copy the same parent cell. In the second case,
the common side is inner and nontrivial, so properness gives equality of
its signals. Induction along the finite path proves the first assertion.

Let $B$ be the small noncentral tile sharing the final print, and choose a
legal parent $C$ of $B$. Its parent form is $\alpha$, so the same form-based
print and path maps are used for $C$ and $D$. The predecessor of the print
on $P_j(d)$ is an ordinary side of $B$, by the routing certificate. Its
third cell equals $C_j(d)$ by the child clause and equals $D_j(d)$ by the
first assertion. The print clause therefore supplies $D_j(d)$ for each
$j\le4$. This argument uses the predecessor in the small owner, rather
than an unsupported claim that the central tile is already legal.
\end{proof}

\begin{lemma}[Exact repaired decomposition]
\label{lem:comp/exact}
\lean{Tilings.Planned.macroblock_is_decomposition,Tilings.Planned.Strict.strict_macroblock_is_decomposition}
\uses{lem:comp/reconstruction,lem:comp/signals,def:control/substitution}
The finite patch $M$ is an exact member of $\widehat\sigma D$.
\end{lemma}
\begin{proof}
Check the clauses cell by cell. Child types and fixed first, second, and
fifth cells agree because each noncentral child has a parent of form $\alpha$.
Remaining inner nonprint fourth cells are $\bot$. Boundary fourth cells carry
$h=H(D)$. For an outer child side $b\subset\sigma'd$, its legal owner $B$
has a parent $C$ with $H(C)=h=H(D)$. Hence its choice between zero and its
fixed positive rank is exactly the choice prescribed by $D$. This is the
case missing from the unchanged composition argument in the source article:
the outer fourth value depends on the parent's mode, even though it is
not an ordinary copied cell.

Lemma~\ref{lem:comp/signals} verifies all ordinary third signals and all
four print cells in the noncentral owners. Remaining third cells are zero
by the child clause. Thus every noncentral tile has the required colours
as a child of $D$. (The outer-side owners used above have boundary sides
by Lemma~\ref{lem:comp/outer-boundary}.) Every side of the central tile is
a complete shared inner side with one of these legal noncentral neighbours. The neighbour forces
equality of the full colour by properness. The repaired substitution
prescribes equal colours to both owners of such a side, including the
borrowed tuple on a print. Consequently the central tile also has all the
required child cells. No initial legality assumption about the centre
has been used.
\end{proof}

\begin{lemma}[Cyclic reconstructed parents are legal]
\label{lem:comp/cyclic-legal}
\lean{Tilings.Planned.reconstructed_cyclic_legal,Tilings.Planned.Strict.strict_reconstructed_cyclic_legal}
\uses{lem:comp/exact,def:certificate,def:control/legal,lem:finite-idempotent}
Suppose the centre $A$ of $M$ is legal and the type $t$ of $D$ is cyclic:
it is the central child type in the template of its own form. Then $D=A$
as typed decorations, and $D$ is legal.
\end{lemma}
\begin{proof}
Both tiles have type $t$. The retained large form is fixed by the central
form map, obtained using Lemma~\ref{lem:finite-idempotent} in the geometric
construction. Its print map is the identity on marked sides. Hence every
marked first-four cell of $A$, as central child of $D$, copies the
corresponding cell of $D$. The remaining cells are local and agree because
both tiles are normal. The header is computed from the type and fourth
vector and therefore agrees as well. Thus $A=D$, proving the claim.
\end{proof}

\begin{lemma}[Noncentral reconstructed parents are legal]
\label{lem:comp/noncentral-legal}
\lean{Tilings.Planned.reconstructed_noncentral_legal,Tilings.Planned.Strict.strict_reconstructed_noncentral_legal}
\uses{lem:comp/exact,lem:comp/cyclic-legal,def:certificate,def:control/legal}
If the centre $A$ is legal and $t$ is noncentral, then $D$ is legal.
\end{lemma}
\begin{proof}
A normal noncentral tile has a marked ordinary inner side with a nonzero
first identity cell, by the separating certificate. Its print in $A$
contains that nonzero value. A legality chain of $A$ with only central
types would copy this marked first cell forever: all parent forms would
be retained large forms, whose print maps are identities. The cell would
have no origin and would be zero, a contradiction. Therefore the chain
reaches a legal noncentral ancestor.

Choose a shortest central-child chain ending at such an ancestor. Passing
to the central child fixes large forms. All but possibly the last two
vertices are therefore cyclic tiles of the same type and, by the argument
of Lemma~\ref{lem:comp/cyclic-legal}, identical decorations. Delete these
repetitions. The chain has length $n\le2$. Length two is impossible here:
the immediate parent of $A$ would have a large form, whereas its known
parent $D$ has a noncentral form. A child's position determines its parent
form, so these cannot both be parents. Thus $A$ has a legal noncentral
parent $E$ immediately above it.

The parents $E,D$ have the same form. Their print maps depend only on this
form, so their marked first-four fields agree through the common child $A$.
Let $\Phi$ be the multiset of type components of first-cell interface identitys
on marked sides. We have $\Phi(E)=\Phi(D)$. The noncentral fingerprint
certificate implies that their types agree. Their local fields then agree
by normality, while their marked first-four fields already agree. Hence
$E=D$, and $D$ is legal.

The fingerprint condition can be checked from exact finite data: original
types have a nonempty inner identity set; a template's connected dual graph
has more than two vertices; small types have enough ordinary marked inner
sides to separate their cardinalities from original types; and each pair
of small types has at least three private ordinary original-neighbour
occurrences. The exceptional special side contributes at most two arbitrary
components. These are finite separating conditions, not informal estimates
that sufficiently large quasi-squares will automatically have them.
\end{proof}

\begin{lemma}[Positive ranks align an entire macroside]
\label{lem:comp/alignment}
\lean{Tilings.Planned.alignment_complete}
\uses{def:certificate,def:tiling,def:control/substitution}
Let the macroblocks of a proper plane tiling of legal typed tiles form a
plane partition with disjoint interiors (legality makes every outer fourth
cell zero or the generated rank of its side; without this a constant
colouring satisfies the hypotheses below on the three-block counterexample).
If every child side of one macroside has a complete opposing side with the
same positive alignment rank, \emph{and every macroside of an opposing
macroblock is mode-uniform} (either all of its child sides carry positive
ranks, or all carry zero), then this macroside has a complete opposing
macroside with the same ordered subdivision.

Without the mode-uniformity hypothesis the statement is false: let $d$
have children of ranks $1,2,3$, each completely matched by the children of
ranks $1,2,3$ of an opposing macroside $e$ which has five children, the
last two facing trivial children of a third block. The hypothesis on $d$
holds, yet $d$ has no complete opposing macroside. This example is a
counterexample only under the weaker legality ``every outer fourth cell is
zero or the generated rank of its side''; under the full-header legality of
the repaired child relation every macroblock of a proper legal tiling has a
common boundary header, so mode-uniformity follows from legality and the
hypothesis is automatic. In
Lemma~\ref{lem:comp/proper} mode-uniformity is available, because every
block is an exact decomposition of its reconstructed parent and the parent
header fixes one mode per side.
\end{lemma}
\begin{proof}
Use the global orientation of ranks along straight sides, required by (G3).
The child side
of rank one coincides completely with a child side of rank one on an
opposing macroblock. Their first endpoints, and thus the beginnings of
their macrosides, coincide. Continue along the finite ordered list of child
segments. An opposing macroblock cannot end inside one of these segments:
the segment has a complete matching side. If it ends between consecutive
segments, the next opposing macroside would have to begin there and its
first rank would be one, whereas the next given rank is larger than one.

The other end needs the mode-uniformity hypothesis. Let the given side
have $L$ children and the opposing side $e$ have $L'$ children, and suppose
$L'>L$. Since child $1$ of $e$ carries a positive rank, mode-uniformity
makes child $L+1$ of $e$ ranked, hence nontrivial, so it needs a complete
match of rank $L+1$. Beyond the end of the given side, on the same side of
the line, lies either the next listed side of the same block, whose first
child has rank $1$, or a side of a different block, which would have to
begin $L$ children earlier and overlap the given block, or a trivial side,
which cannot match a nontrivial one. All three are impossible, so $L'=L$.

For completeness, an alternative opposing macroside cannot have begun
earlier and merely join the contact at this switch. An open portion of that
earlier side would lie opposite the same open portion of the original
opposing block, with their interiors on the same half-plane side. Local
polygonal interior neighbourhoods would overlap, contradicting the
disjointness of macroblock interiors. Coverage supplies an opposing block
after each endpoint, so the rank-reset obstruction cannot be avoided by
leaving a gap. Applying the argument at the other end shows that the
opposing side terminates exactly with the given side. Complete child-side
matching then identifies the whole ordered subdivision. This geometric
lemma must be formalized separately from the algebraic header kernel.
\end{proof}

\begin{lemma}[The reconstructed plane tiling is proper]
\label{lem:comp/proper}
\lean{Tilings.Planned.composed_proper,Tilings.Planned.Strict.strict_composed_proper}
\uses{lem:comp/assembly,lem:comp/exact,lem:comp/alignment,lem:comp/reconstruction,lem:comp/boundary-header}
For a proper legal plane tiling $\mathcal T$, replace each assembled
macroblock by its reconstructed parent and scale down by the substitution
factor. Each assembled macroblock is replaced by its reconstructed parent in
the sense of Lemma~\ref{lem:comp/reconstruction}, i.e.\ the parent whose
type is the type of the common boundary header of the block given by
Lemma~\ref{lem:comp/boundary-header}. The result $\mathcal T'$ is a proper tiling of normal identity-correct decorations
with active inner sides, and $\mathcal T$ is an exact repaired decomposition
of $\mathcal T'$.
\end{lemma}
\begin{proof}
Assembly partitions the plane into the macroblocks. Their underlying
polygons are scaled parent polygons, so replacing and scaling them preserves
coverage and disjoint interiors. Exact decomposition follows block by block
from Lemma~\ref{lem:comp/exact}; normality, identity-correctness, and activity of inner sides follow
from Lemma~\ref{lem:comp/reconstruction}.

Every macroside of every block is mode-uniform: the block is an exact
decomposition of its reconstructed parent $E$ (Lemma~\ref{lem:comp/exact}),
and the outer fourth cells on $\sigma'e$ are all zero or all positive
ranks according to $z_E(e)$. So the alignment lemma applies.

Consider a nontrivial side $d$ of a reconstructed parent $D$. If it is
inner, $z_D(d)=0$ by definition. If it is outer, its fifth cell is zero,
so nontriviality gives $D_4(d)\ne0$ and again $z_D(d)=0$. Thus all its
child outer fourth cells are positive ranks. Properness of $\mathcal T$
matches each completely with an equal positive rank. The alignment lemma
gives a complete opposing parent side and the same subdivision. Its five
ports consequently coincide with the corresponding five ports on the
given side. The third cells on those ports agree in $\mathcal T$. These
five values are exactly the five reconstructed parent cells, proving full
colour equality. A trivial outer side imposes no requirement. Hence
$\mathcal T'$ is proper.
\end{proof}

\begin{lemma}[Central acyclic reconstructed parents are legal]
\label{lem:comp/acyclic-legal}
\lean{Tilings.Planned.reconstructed_acyclic_legal,Tilings.Planned.Strict.strict_reconstructed_acyclic_legal}
\uses{lem:comp/proper,lem:comp/exact,lem:comp/noncentral-legal,lem:comp/cyclic-legal,def:certificate}
Every reconstructed parent of central acyclic type in $\mathcal T'$ is legal.
\end{lemma}
\begin{proof}
Let $D$ be such a parent, of type $t$, and let $A$ be its original legal
central child. Every noncentral or cyclic tile of $\mathcal T'$ is already
legal by the preceding legality lemmas. All reconstructed tiles are normal
and identity-correct with active inner sides, so the assembly lemma also applies to
$\mathcal T'$. In the second-level macroblock containing $D$, every
noncentral tile is legal. The full repaired reconstruction and exactness
lemmas give a normal parent $C$ of type $r$, with $D$ its central child.
No unproved inheritance of a single-scale flag rule is needed.

The type $r$ is noncentral. If it were central, its large form would be
fixed by the central-child form map. It and its central child $t$ would
have the same large form, and hence the same central child type, making
$t$ cyclic. This contradicts the assumption on $D$.

The nonzero marked first cell of normal noncentral $C$ is copied into
$D$ and then into $A$. The origin argument from the noncentral legality
lemma therefore shows that a legal ancestry chain of $A$ reaches a
noncentral type; removing cyclic repetitions gives length at most two.
An immediate parent of $A$ has the same large form as $D$, and hence a
central type. The chain must have length exactly two:
\[
 A\ \longleftarrow\ E\ \longleftarrow\ F,
\]
where $E,F$ are legal and $F$ is noncentral.

The parents $E,D$ have the same form and agree in their marked first-four
fields through $A$, so $\Phi(E)=\Phi(D)$. Write $L=4N$ for the number
of marked sides of a large form, and $Q$ for the maximum on a noncentral
form. At most $Q$ marked sides of $D$, the central child of $C$, are prints.
On each of the remaining at least $L-Q$ sides its first interface identity
contains $t$. Thus the multiplicity of $t$ in $\Phi(D)$ is at least $L-Q$.
Suppose $E$ had a different central type $t'$. Its nonprint identities
contain $t'$ and a small neighbour type, neither equal to $t$. Only its
at most $Q$ prints can contribute $t$; a copied print cell contributes at
most one occurrence of $t$ per print under the type projection, so the
multiplicity is at most $Q$, and in any case at most $2Q$.
This contradicts $L-Q>2Q$, supplied by the certificate. Therefore $E$
also has type $t$. Normal local cells and the marked copied cells now
identify $E=D$. Since $E$ is legal, $D$ is legal.
\end{proof}

\begin{theorem}[The repaired composition theorem]
\label{thm:compose}
\lean{Tilings.Planned.compose_legal_exists}
\uses{def:certificate,def:hereditary,lem:comp/proper,lem:comp/noncentral-legal,lem:comp/cyclic-legal,lem:comp/acyclic-legal}
For every macrotemplate system satisfying Definition~\ref{def:certificate}
with hereditary listed sides (Definition~\ref{def:hereditary})
and the explicit header substitution of Definition~\ref{def:control/substitution},
every proper legal plane tiling has a proper exact repaired composition
whose tiles are legal.
\end{theorem}
\begin{proof}
Construct $\mathcal T'$ by assembly and reconstruction. The properness
lemma proves that it is a plane tiling, is proper, and decomposes exactly
to the input. Its reconstructed types are either noncentral, central cyclic,
or central acyclic. The three legality lemmas cover these disjoint cases.
Thus every tile of $\mathcal T'$ is legal. The conclusion uses the same
properness and legality predicates as the input, so the theorem can be
iterated. The header modification is what provides this closure.
\end{proof}

\begin{theorem}[Lifting an actual side-to-side-supertiling hierarchy]
\label{lem:lift-hierarchy}
\lean{Tilings.Planned.lift_hierarchy}
\uses{def:certificate,def:hereditary,def:hierarchy,def:control/origin,def:control/substitution}
Suppose the macrotemplate system has hereditary listed sides, so that no
fine outer side straddles a parent listed vertex and ``the unique
containing side'' below is defined. Suppose an actual $\sigma'$ hierarchy
is given and its finite supertiles
are side-to-side at every level: for every tile $X$ of level $n$ and every
$k\le n$ the patch $\sigma'^{\,k}X$ of level-$(n-k)$ tiles is side-to-side
(the full decomposition alone is not enough: two level-$1$ siblings with
partially overlapping sides would give coincident fine sides different port
ranks). Then it admits compatible legal decorations at every
level under the repaired substitution; its decorated level-zero tiling
is proper. This side-to-side-supertiling assumption is an additional
geometric instantiation obligation, supplied by a correctly grouped STS
construction, not proved by the header kernel.
\end{theorem}
\begin{proof}
At each level the hierarchy specifies each child's type and its parent.
On an outer side, follow the unique containing side through parent levels.
If it stays outer forever, put its fourth cell equal to zero at all levels
of that chain. If it eventually becomes inner, assign the fixed positive
alignment ranks on the outer levels below the first inner side. These
choices obey the outer-fourth substitution clause. Set remaining inner nonprint
fourth cells to $\bot$, and boundary fourth cells to the computed header
of their parent.

Set every remaining noncopied fresh cell to its prescribed child value;
this includes fixed local cells and first identities on special nonprints.
For a recursively copied cell, follow its source: use the origin value if an origin is reached, and
zero if it is not. Copy-source uniqueness makes this assignment
unambiguous. Equality along a copy edge holds in both cases. A header depends
only on the type and the outer fourth cells, all of which were assigned
before the copying step: inner entries of its mode vector are false regardless
of their fourth values. Its assignment therefore introduces no additional
recursive choice. The child clauses are now
all satisfied. Each actual ancestor chain witnesses legality, including
the no-origin condition at every level.

Inside a macroblock, shared side colours agree because both owners receive
the same parent data; synthesized boundary headers also agree. For an outer
side with nonzero fourth cell, its ancestor side becomes inner at a finite
level. Both overlapping sides therefore belong to one finite higher
supertile. The assumed STS property makes their descendant sides coincide
completely, with matching ranks and ports. Their copied cells either reach
the same shared origin, or matching origins on a shared side, or both have
no origin and are zero. The second alternative is an induction on the
level: the two copied cells on the coincident sides $b\subset\sigma'd$ and
$b'\subset\sigma'd'$ copy the cells of the parent sides $d,d'$, which are
again coincident sides in the next level with one fewer level to go before
the chain becomes inner; the induction hypothesis gives equality there.
The third cell of an outer side is nonzero only when that side is a port,
by the inner-side clause of (G4); hence no outer side other than a port
carries a parent signal, and the coincident side of a port is the port of
the same number of the opposing block. Hence their full colours agree. The same argument
applies to inner sides, already lying in a finite supertile. An outer side
whose chain stays outer forever has zero fourth and fifth cells and is
trivial, so its shifted contacts are permitted. The level-zero decoration
is proper, and the construction is compatible across all hierarchy levels.
\end{proof}

\begin{theorem}[Soundness and the conditional representation]
\label{thm:comp/soundness}
\lean{Tilings.Planned.projects_to_hierarchy}
\uses{thm:compose,lem:lift-hierarchy,def:hierarchy,def:certificate,def:hereditary}
Under the certificate assumptions, with hereditary listed sides, erasing
the colours of every proper
legal tiling gives a $\sigma'$-hierarchical tiling. If in addition the
actual hierarchies to be represented have side-to-side finite supertiles,
then the erasures of proper legal tilings are exactly those hierarchical
tilings. The resulting decorated tile alphabet is finite.
\end{theorem}
\begin{proof}
Apply the repaired composition theorem repeatedly. Classical dependent
choice, or recursive selection of a witness for each natural-number level,
gives proper legal tilings $\mathcal T^{(0)},\mathcal T^{(1)},\ldots$,
each an exact composition of its predecessor. Erasure commutes with the
repaired substitution, since the underlying geometry was unchanged.
Thus their erased sequence witnesses the required $\sigma'$ hierarchy.
This direction uses no assumed existence of a coloured hierarchy.

Conversely, the lifting theorem supplies a proper legal decoration of
every hierarchy satisfying its geometric STS condition. Finiteness follows
because there are finitely many forms, types, side positions, and ordinary
index tags, while the dependent header sum is finite. Legality is a subset
of this finite alphabet. These two implications prove the conditional
representation statement.

This theorem is deliberately conditional on the geometric certificate and
the stated realization condition. It is not a completed proof of the
general STS/FLC soficity theorems. Establishing a valid connected grouping,
the routing and separation certificates, the compatibility with the
original substitution, and the repaired FLC geometric reduction remains
outside this chapter. In the crown-bearing variants, positive outer
fourth cells make every header mode vector false; the header then carries
the former parent type with redundant control bits. Crown consistency and
the additional language argument require their own nodes and are not
inferred from the conditional hierarchical representation proved here.
\end{proof}
